\section{Análise de Sensibilidade das Instâncias e Instâncias Adversariais}
\label{sec:sensibilidade}

Para avaliar como a geometria e a topologia das instâncias afetam a qualidade das soluções heurísticas e aproximadas (e o esforço de busca do \textit{Branch-and-Bound}), extraímos um conjunto de \textit{features} estruturais para todas as 14 instâncias da TSPLIB e correlacionamos essas métricas com os \textit{gaps} empíricos em relação ao ótimo conhecido.

\subsection{Extração de Características Geométricas e Topológicas}

A partir das coordenadas 2D de cada instância (\texttt{NODE\_COORD\_SECTION}, \texttt{DISPLAY\_DATA\_SECTION} em \texttt{bays29}/\texttt{dantzig42} e projeção via \textit{Classical MDS} nas 5 instâncias puramente \texttt{EXPLICIT}) e da matriz exata de distâncias, extraíram-se as seguintes \textit{features} (salvas em \texttt{results/features\_instancias\_sensibilidade.csv}):

\begin{itemize}
    \item \textbf{Razão de Aspecto (\textit{Aspect Ratio}, $\mathrm{AR}$):} razão entre o maior e o menor lado do \textit{bounding box} dos pontos ($\max(\Delta x, \Delta y) / \min(\Delta x, \Delta y)$), medindo o alongamento geométrico da instância.
    \item \textbf{Proporção no Fecho Convexo ($\mathrm{Hull}$):} fração de vértices pertencentes ao fecho convexo 2D ($|H|/N$, calculado via \textit{Monotone Chain} de Andrew).
    \item \textbf{Entropia Espacial ($H_{\text{esp}}$):} entropia de Shannon normalizada em $[0,1]$ sobre uma grade $4 \times 4$ no \textit{bounding box}. Valores próximos de $1$ indicam distribuição espacial uniforme; valores baixos indicam forte \textit{clustering} (aglomeração).
    \item \textbf{Estatísticas da MST ($\Delta_{\text{MST}}$ e $|V'|/N$):} grau máximo de um vértice na Árvore Geradora Mínima (presença de \textit{hubs}) e proporção de vértices de grau ímpar.
    \item \textbf{Coeficiente de Variação ao Vizinho Mais Próximo ($\mathrm{CV}_{\text{NN}}$):} razão $\sigma/\mu$ das distâncias de cada vértice ao seu vizinho mais próximo, capturando heterogeneidade local e presença de \textit{outliers}.
\end{itemize}

A Tabela~\ref{tab:features_instancias} consolida as \textit{features} extraídas e os \textit{gaps} percentuais de cada método (\textit{Convex Hull}, \textit{Christofides}, \textit{Nearest Neighbor}, \textit{Twice-Around-the-Tree} e folga da 1-árvore na raiz $\mathrm{Gap}_{LB_0}$).

\begin{table}[htbp]
\centering
\caption{Características geométricas/topológicas e \textit{gaps} empíricos (\%) das 14 instâncias da TSPLIB.}
\label{tab:features_instancias}
\resizebox{\columnwidth}{!}{%
\begin{tabular}{l|c|r|r|r|c|r||r|r|r|r|r}
\hline
\textbf{Instância} & \textbf{$N$} & \textbf{AR} & \textbf{Hull} & \textbf{$H_{\text{esp}}$} & \textbf{$\Delta_{\text{MST}}$} & \textbf{$\mathrm{CV}_{\text{NN}}$} & \textbf{CH (\%)} & \textbf{Chr (\%)} & \textbf{NN (\%)} & \textbf{TAT (\%)} & \textbf{$\mathrm{Gap}_{LB_0}$} \\ \hline
\texttt{burma14}   & 14 & 2.00 & 35.7\% & 0.84 & 4 & 0.76 & 0.0  & 8.5  & 21.8 & 22.6 & 23.5 \\
\texttt{ulysses16} & 16 & 4.05 & 43.8\% & 0.71 & 3 & 1.31 & 1.6  & 2.5  & 45.6 & 12.3 & 30.8 \\
\texttt{gr17}      & 17 & 2.39 & 35.3\% & 0.71 & 3 & 0.74 & 0.0  & 5.4  & 4.9  & 12.8 & 28.0 \\
\texttt{gr21}      & 21 & 1.17 & 38.1\% & 0.70 & 3 & 0.65 & 2.0  & 17.0 & 23.1 & 37.7 & 16.8 \\
\texttt{ulysses22} & 22 & 4.05 & 31.8\% & 0.72 & 3 & 1.63 & 0.4  & 4.6  & 51.0 & 13.5 & 30.6 \\
\texttt{gr24}      & 24 & 1.41 & 29.2\% & 0.85 & 4 & 0.32 & 3.5  & 14.4 & 22.1 & 31.1 & 15.0 \\
\texttt{fri26}     & 26 & 1.55 & 34.6\% & 0.83 & 3 & 0.67 & 0.3  & 10.9 & 18.7 & 29.3 & 12.1 \\
\texttt{bays29}    & 29 & 1.00 & 24.1\% & 0.90 & 4 & 0.32 & 2.1  & 16.3 & 11.8 & 38.3 & 19.7 \\
\texttt{dantzig42} & 42 & 1.84 & 19.0\% & 0.96 & 4 & 0.60 & 3.6  & 15.0 & 36.8 & 42.2 & 14.2 \\
\texttt{swiss42}   & 42 & 1.68 & 21.4\% & 0.89 & 3 & 0.45 & 5.3  & 8.2  & 28.0 & 33.5 & 13.0 \\
\texttt{att48}     & 48 & 1.50 & 22.9\% & 0.86 & 4 & 0.50 & 2.2  & 18.4 & 21.0 & 30.6 & 15.0 \\
\texttt{eil51}     & 51 & 1.09 & 15.7\% & 0.99 & 4 & 0.26 & 3.8  & 13.6 & 20.0 & 45.8 & 9.6  \\
\texttt{berlin52}  & 52 & 1.47 & 15.4\% & 0.81 & 3 & 0.78 & 7.5  & 14.0 & 19.1 & 37.2 & 18.2 \\
\texttt{st70}      & 70 & 1.06 & 14.3\% & 0.98 & 4 & 0.49 & 3.9  & 15.0 & 23.0 & 19.4 & 15.0 \\ \hline
\end{tabular}%
}
\end{table}

\subsection{Modelagem Preditiva e Análise de Instâncias Adversariais}

Para quantificar o peso de cada característica no desempenho dos algoritmos, ajustaram-se modelos de \textbf{Regressão Linear Múltipla} com \textit{features} padronizadas (via \texttt{StandardScaler}, de modo que os coeficientes $\beta$ expressem o impacto em pontos percentuais por desvio-padrão) e \textbf{Árvores de Decisão de Regressão} (\texttt{DecisionTreeRegressor}, profundidade máxima $2$ para evitar sobreajuste com $14$ amostras). A Tabela~\ref{tab:modelos_sensibilidade} e a Figura~\ref{fig:sensibilidade_instancias} resumem os resultados.

\begin{table}[htbp]
\centering
\caption{Desempenho dos modelos de sensibilidade (Regressão Múltipla e Árvore de Decisão, 5 \textit{features}).}
\label{tab:modelos_sensibilidade}
\resizebox{\columnwidth}{!}{%
\begin{tabular}{l|c|c|c|c|l}
\hline
\textbf{Alvo (Gap \%)} & \textbf{Reg.\ $R^2$} & \textbf{Reg.\ RMSE} & \textbf{Árv.\ $R^2$} & \textbf{Árv.\ RMSE} & \textbf{Fatores Dominantes (Correlação $r$ e Peso $\beta$)} \\ \hline
\textit{Convex Hull}      & $0.6827$ & $1.18$ p.p. & --       & --          & $\mathrm{Hull}$ ($r=-0.74$, $\beta=-2.16$) \\
\textit{Christofides}     & $0.8323$ & $1.99$ p.p. & $0.8506$ & $1.88$ p.p. & $\mathrm{AR}$ ($r=-0.84$, $\beta=-5.66$); $\Delta_{\text{MST}}$ ($r=+0.57$, $\beta=+2.38$) \\
\textit{Nearest Neighbor} & $0.7036$ & $6.45$ p.p. & --       & --          & $\mathrm{CV}_{\text{NN}}$ ($r=+0.71$, $\beta=+9.04$); $\mathrm{AR}$ ($r=+0.72$, $\beta=+4.06$) \\
\textit{Twice-Around}     & $0.5687$ & $7.08$ p.p. & $0.7540$ & $5.35$ p.p. & $\mathrm{AR}$ ($r=-0.73$, $\beta=-5.67$); $\mathrm{CV}_{\text{NN}}$ ($r=-0.69$) \\
$\mathrm{Gap}_{LB_0}$ (B\&B) & $0.8499$ & $2.58$ p.p. & $0.8873$ & $2.24$ p.p. & $\mathrm{AR}$ ($r=+0.84$, $\beta=+4.20$); $H_{\text{esp}}$ ($r=-0.73$, $\beta=-5.39$) \\ \hline
\end{tabular}%
}
\end{table}

\begin{figure}[htbp]
    \centering
    \includegraphics[width=\columnwidth]{analise_sensibilidade_instancias.png}
    \caption{Análise de sensibilidade topológica: matriz de correlação de Pearson entre \textit{features} e \textit{gaps}, pesos padronizados da regressão múltipla e mapa de instâncias adversariais com o corte principal da árvore de decisão ($\mathrm{AR}=1{,}92$).}
    \label{fig:sensibilidade_instancias}
\end{figure}

A análise conjunta das correlações, regressões e regras da árvore de decisão revela um fenômeno central de \textbf{adversarialidade oposta entre famílias de algoritmos}:

\begin{enumerate}
    \item \textbf{Instâncias Alongadas com \textit{Outliers} (\texttt{ulysses16}, \texttt{ulysses22}, \texttt{gr17}) -- Adversariais para \textit{Nearest Neighbor} e B\&B, Favoráveis para \textit{Christofides}/\textit{Twice-Around}:} Com razão de aspecto elevada ($\mathrm{AR} \in [2{,}39;\, 4{,}05]$), baixa entropia ($H_{\text{esp}} \approx 0{,}71$) e forte heterogeneidade local ($\mathrm{CV}_{\text{NN}}$ até $1{,}63$ em \texttt{ulysses22}), essas instâncias provocam os piores \textit{gaps} do \textit{Nearest Neighbor} ($+45{,}6\%$ em \texttt{ulysses16} e $+51{,}0\%$ em \texttt{ulysses22}, pois a escolha míope visita aglomerados densos e sofre saltos longos no fechamento) e abrem a maior folga no limite inferior de 1-árvore do B\&B ($\mathrm{Gap}_{LB_0} \approx 30{,}7\%$, regra da árvore para $\mathrm{AR} > 3{,}22$, explicando o colapso precoce de \texttt{ulysses22} com $N=22$). Em contrapartida, nessa geometria quase unidimensional a MST aproxima-se de um caminho simples ($\Delta_{\text{MST}}=3$, poucos vértices ímpares), tornando \textit{Christofides} ($2{,}5\%$ a $4{,}6\%$) e \textit{Twice-Around} ($12{,}3\%$ a $13{,}5\%$) altamente precisos.
    \item \textbf{Instâncias Compactas, Uniformes e com \textit{Hubs} na MST (\texttt{eil51}, \texttt{bays29}, \texttt{dantzig42}, \texttt{att48}, \texttt{gr21}) -- Adversariais para Algoritmos Baseados em MST:} A árvore de decisão particiona nitidamente o conjunto em $\mathrm{AR} \le 1{,}92$: instâncias isotrópicas ($\mathrm{AR} \approx 1{,}00$--$1{,}50$) com alta entropia espacial ($H_{\text{esp}} \in [0{,}86;\, 0{,}99]$) e baixa variação local ($\mathrm{CV}_{\text{NN}} \le 0{,}32$) elevam o erro médio do \textit{Twice-Around} para $42{,}0\%$ (pico de $+45{,}8\%$ em \texttt{eil51} e $+42{,}2\%$ em \texttt{dantzig42}) e do \textit{Christofides} para $15{,}5\%$ (pico de $+18{,}4\%$ em \texttt{att48}). Isso decorre da formação de \textit{hubs} de grau $4$ na MST ($r=+0{,}57$ com \textit{Christofides}) e alta fração de vértices ímpares ($|V'|/N = 62{,}5\%$ em \texttt{att48}, $r=+0{,}59$, Spearman $\rho=+0{,}67$), ampliando o custo do emparelhamento perfeito e gerando desvios mais custosos na etapa de \textit{shortcutting}.
    \item \textbf{Sensibilidade do \textit{Convex Hull} à Proporção de Pontos na Fronteira:} A heurística de inserção por fecho convexo exibiu forte correlação negativa com $\mathrm{Hull}$ ($r=-0{,}74$, Spearman $\rho=-0{,}85$, $\beta=-2{,}16$ com $R^2=0{,}6827$). Quando mais de $35\%$ dos pontos já compõem a envoltória externa (\texttt{burma14}, \texttt{gr17}), a heurística atinge o ótimo global ($0{,}0\%$ de erro); inversamente, em \texttt{berlin52} (apenas $15{,}4\%$ no fecho e alta densidade interna com $\mathrm{CV}_{\text{NN}}=0{,}78$), o grande número de inserções gulosas sequenciais eleva o erro ao máximo observado de $+7{,}5\%$.
\end{enumerate}
